\section{等效宽度（Effective Width）}

\subsection{从 FLOPs 角度理解}

\begin{importantbox}{等效宽度的定义}
MoE 的等效宽度 = 激活 Expert 数 $\times$ 每个 Expert 的中间维度：
$$
d_{\text{ff}}^{\text{eff}} = k \times m = 8 \times 768 = 6144 = 3 \times d_{\text{model}}
$$
而 Dense 模型的 FFN 宽度为 $25600 = 5 \times d_{\text{model}}$。
\end{importantbox}

两种架构的 FLOPs 计算：

\textbf{Dense FFN}（含 SwiGLU 的三个矩阵 up/gate/down）：
$$
\text{FLOPs}_{\text{Dense}} = 3 \times d_{\text{model}} \times d_{\text{ff}}
$$

\textbf{MoE}（$k$ 个 Expert 各执行一次 up/gate/down）：
$$
\text{FLOPs}_{\text{MoE}} = k \times 3 \times d_{\text{model}} \times m = 3 \times d_{\text{model}} \times (k \cdot m)
$$

令两者 FLOPs 相等：$d_{\text{ff}} = k \cdot m$，这就是等效宽度。

\subsection{加权求和 vs. Concat：从 Rank 角度理解}

\begin{warningbox}{MoE 输出不是 Concat 而是加权求和}
8 个 Expert 输出的不是拼接成 $8 \times 768 = 6144$ 维的向量，而是在 $d_{\text{model}} = 2048$ 的空间里做\textbf{加权求和}，最终输出仍为 2048 维。
\end{warningbox}

为什么加权求和仍能等效于 6144 的宽度？有两个视角：

\textbf{视角 1：FLOPs 等价}。无论是否 concat，参与实际计算的矩阵乘法总量相同。

\textbf{视角 2：矩阵 Rank 等价}。加权求和可以改写为：

$$
\sum_{i=1}^{k} \alpha_i \cdot v_i = \begin{bmatrix} \alpha_1 I_d & \alpha_2 I_d & \cdots & \alpha_k I_d \end{bmatrix} \begin{bmatrix} v_1 \\ v_2 \\ \vdots \\ v_k \end{bmatrix}
$$

其中 $I_d$ 是 $d \times d$ 的单位矩阵。这等价于一个 $d \times kd$ 的块对角稀疏矩阵乘以一个 $kd$ 维的 concat 向量。从矩阵 rank 角度看，表达能力（多样性）是等价的：
\begin{itemize}
  \item 单个 Expert 的表达能力被降低（维度变瘦）
  \item 多个 Expert 共同承载了足够的多样性
\end{itemize}

\subsection{本章小结}

等效宽度 $k \times m$ 是理解 MoE 计算量与表达能力的关键概念。Qwen3-30B-A3B 的等效宽度为 $3 \times d_{\text{model}}$，低于 Dense 的 $5 \times d_{\text{model}}$，但通过 128 个 Expert 的总参数量保持了巨大的模型容量。

